\begin{lemma}[型フィルタは代入可能性を保つ]
\label{lem:type-filter}
$L$ が $(f,h)$-タプル代入可能で $F\subseteq M$ なら，
$L\cap h^{-1}(F)$ も $(f,h)$-タプル代入可能である．
\end{lemma}

\begin{proof}
$\tuple x,\tuple y$ が等しい $h$-ベクトルをもち，
$L':=L\cap h^{-1}(F)$ において文脈 $E$ を共有するとする．
両者は $L$ でも $E$ を共有するので，$L$ における分布が等しい．
$E'[\tuple x]\in L'$ なら $E'[\tuple y]\in L$ であり，
$h(E'[\tuple y])=h(E'[\tuple x])\in F$ なので $E'[\tuple y]\in L'$ である．
\end{proof}

\begin{corollary}[pseudovariety の二分法]
\label{cor:pseudovariety-dichotomy}
$\mathbf V$ を有限半群の pseudovariety とする．
\begin{enumerate}[label=\textup{(\roman*)},leftmargin=*]
\item $\mathbf V\subseteq\mathbf{LG}$ なら，任意のアルファベット，
$h(\Sigma^+)\in\mathbf V$ となる任意の型付け $h$，任意の固定された $f,r_{\max}$ について，
定理~\ref{thm:local-group-polydata} のとおり総長 $O(\|G\|^3\bar\theta_G)$ の特徴データが存在する．
\item $\mathbf V\not\subseteq\mathbf{LG}$ なら $U_1\in\mathbf V$ である．
ファンアウト1において，$h_{\mathrm U}(\Sigma^+)=U_1$ となる型付け $h_{\mathrm U}$ と，
$\|G\|,\theta_G=O(n)$ の標的が存在し，その標的の
$\mathcal B^{\le2}_{1,h_{\mathrm U}}$ に対する特徴標本はどれも長さ $3\cdot2^n-2$ の語を含む．
$\Cmcf{1}{h_{\mathrm U}}$ 全体で特徴標本をもつ任意の集合駆動作用素についても同じである．
\end{enumerate}
\end{corollary}

\begin{proof}
(i) は定理~\ref{thm:local-group-polydata} である．
(ii) について，$S\in\mathbf V\setminus\mathbf{LG}$ を選ぶ．
定理~\ref{thm:lg-characterization} により $S$ は $U_1$ と同型な部分半群をもち，
pseudovariety は部分半群で閉じている．

\cite[Proposition~6.3]{kuriyamaCFG} の倍化族 $L_n$，型付け $h_{\mathrm U}$，
型1の長い語 $w_n^*$ を取る．
そこでは $L_n$ が Clark--Eyraud 代入可能であることが示されており，
したがって $(1,h_{\mathrm U})$-タプル代入可能である．
補題~\ref{lem:type-filter} により
$L_n^-:=L_n\cap h_{\mathrm U}^{-1}(\zeta)=L_n\setminus\{w_n^*\}$ もそうである．
$K\subseteq L_n$ が $w_n^*$ を含まなければ $K\subseteq L_n^-$ であり，
命題~\ref{prop:reconstruction-soundness} により
$L(\mathcal B^{\le2}_{1,h_{\mathrm U}}(K))\subseteq L_n^-$ となる．
よってどの特徴標本も $w_n^*$ を含む．
集合駆動作用素についての主張は，対 $L_n,L_n^-$ に対する
\cite[Proposition~6.3(iv)]{kuriyamaCFG} の区別対論法である．
\end{proof}

\begin{remark}[二分法の適用範囲]
\label{rem:dichotomy-scope}
\begin{enumerate}[label=\textup{(\roman*)},leftmargin=*]
\item (i) はすべての固定ファンアウトで成り立つが，(ii) はファンアウト1でしか示されていない．
倍化族はそのままでは持ち上がらない．
実際 $L_2$ は $(2,h_{\mathrm U})$-タプル代入可能でない．
型1の対 $(\mathtt[,\mathtt a\mathtt])$ と $(\mathtt[\mathtt a,\mathtt])$ は
文脈 $\mathtt{[[aa]}\,\blank_1\,\mathtt a\,\blank_2\,\mathtt]$ を共有するが，分布が異なる．
これは $n\le3$ について網羅的に確かめた．
各成分の文字集合を記録する，より細かい型付け（像はやはり半束）でも，$n=3$ で同様の破綻が現れる．
(ii) がファンアウト $f\ge2$ でも成り立つかは，
未解決問題~\ref{op:fanout-two} の pseudovariety 版である．
\item この二分法は pseudovariety についてのものである．
$h(\Sigma^+)\notin\mathbf{LG}$ となる個々の型付け $h$ について，
$\Cmcf{1}{h}$ が倍化族の符号化された複製を含むとは限らない．
そのような型付けがすべて超多項式の特徴データを強いるかは未解決である．
\end{enumerate}
\end{remark}
